\begin{exercise}{108}
In this exercise, we will use minimal fibrations to show that the geometric
realization of a Kan fibration $p\colon X\to Y$ is a Serre fibration, i.e., it has the RLP
with respect to the inclusions $D^n\times\{0\}\to D^n\times D^1$. Here, $D^n$ is the
$n$-dimensional topological cube.
\begin{enumerate}[label=(\arabic*)]
 \item Show that the geometric realization of a minimal fibration is a Serre fibration whose fibre
       is given by the geometric realization of the fibre of the Kan fibration.
 \item Show that the geometric realization of a trivial fibration is a Serre fibration.
\end{enumerate}
Hints: For (1), show that the realization of a locally trivial map is also locally trivial. Then
show that a locally trivial map of spaces is a Serre fibration. For (2), show that there exists
a monomorphism $X\to W$ with $W$ a contractible Kan complex. Consider the maps
$X\to W\times Y\to Y$. Show that the latter is a trivial fibration and deduce that $p$ is a
retract of $W\times Y\to Y$.

\begin{proof}
Use compactly generated spaces; products with the topological cubes have
their ordinary topology. Realization preserves products by
Remark~1.2.60 and preserves colimits by its adjunction with $\mathcal S$.
We first check the pullback comparison needed for the fibre assertion.
By Corollary~1.1.41, simplicial subsets realize as closed CW subcomplexes.
If $A\subseteq B$ and $f\colon T\to B$, then
\[
 |f^{-1}(A)|=|f|^{-1}(|A|)
\]
as subspaces of $|T|$: an interior point of a nondegenerate simplex
$t$ maps into the open cell of the nondegenerate part of $f(t)$,
since collapsing degeneracies keeps all barycentric coordinates positive.
It therefore maps into $|A|$ precisely when $f(t)\in A$.
Apply this to the diagonal $B\hookrightarrow B\times B$ and use
preservation of products. For arbitrary maps $T\to B\leftarrow V$ this gives
\[
 |T\times_B V|\cong |T|\times_{|B|}|V|.
\]

(1) Suppose $p$ is minimal. Put $q=|p|$, $E=|X|$, $B=|Y|$,
$B_m=|\mathrm{sk}_mY|$, and $E_m=q^{-1}(B_m)$.
The assumed simplicial local triviality gives, for each nondegenerate
$m$-simplex $\sigma$ of $Y$, a product trivialization over
$D_\sigma=|\Delta^m|$ with fibre $P_\sigma=|F_\sigma|$.
Write $\chi_\sigma\colon D_\sigma\to B_m$ for its characteristic map and
$\kappa_\sigma\colon D_\sigma\times P_\sigma\to E_m$ for the map on total spaces.
The pullback comparison above implies, on the entire boundary, a homeomorphism
\[
 \partial D_\sigma\times P_\sigma
 \xrightarrow{\cong}
 \partial D_\sigma\times_{B_{m-1}}E_{m-1},
 \qquad (a,z)\longmapsto(a,\kappa_\sigma(a,z)).
\]
In particular this identifies the full old fibre at each boundary point,
including when faces are identified or degenerate.
Pulling back the simplicial skeletal attachments along $p$ and realizing gives
\[
 E_m=E_{m-1}\amalg_{\coprod_\sigma(\partial D_\sigma\times P_\sigma)}
                  \coprod_\sigma(D_\sigma\times P_\sigma),
\]
with the analogous attachment formula for $B_m$.
Here pullback preserves colimits in a slice of simplicial sets, as is
checked degreewise in the fibres of sets.

We construct open trivializing neighborhoods through these attachments.
Start with the open cell containing a chosen point of $B$, say of
dimension $d$, and its product chart with fibre $P$.
Suppose $U\subseteq B_{m-1}$ has a chart
$\psi\colon q^{-1}(U)\cong U\times P$.
For a new cell $\sigma$, put
$V_\sigma=(\chi_\sigma|_{\partial D_\sigma})^{-1}(U)$.
The boundary homeomorphism just established and $\psi$ give a homeomorphism
\[
 V_\sigma\times P_\sigma\cong V_\sigma\times P,
 \qquad (a,z)\longmapsto(a,h_\sigma(a,z)).
\]
Both $h_\sigma$ and its fibrewise inverse are jointly continuous.
Choose a boundary collar $c_\sigma(a,t)=(1-t)a+t b_\sigma$,
$0\le t<1/2$, where $b_\sigma$ is the barycenter of $D_\sigma$.
Adjoin the collars of $V_\sigma$ to $U$ and extend the chart by
\[
 \kappa_\sigma(c_\sigma(a,t),z)
 \longmapsto
 \bigl(\chi_\sigma(c_\sigma(a,t)),h_\sigma(a,z)\bigr).
\]
The new set is open in $B_m$ and meets $B_{m-1}$ in exactly $U$:
its full preimage in every attaching disk is the indicated open collar.
At $t=0$ the formula agrees with $\psi$ by the definition of
$h_\sigma$, even across folded boundary points.
The inverse uses $h_\sigma(a,-)^{-1}$ and has the same compatibility.
Both maps therefore descend continuously through the attachment quotients:
in compactly generated spaces product with $P$ preserves quotient
colimits by the exponential law.

Continue over all higher skeleta. The resulting union of neighborhoods
is open by the weak topology, and the compatible charts and their inverses
give a product chart on its full inverse image.
Indeed, every simplex of $X$ lies over some finite skeleton of $Y$,
so both spaces have the corresponding colimit topologies; products
again preserve these colimits.
This proves local triviality of $q$.
The same pullback comparison identifies its fibre over a vertex $y$
canonically with $|X_y|$, and its fibre over an interior point of
$\sigma$ with $P_\sigma$.

For completeness, local triviality gives the required Serre lifting
property as follows. Given a lifting problem on $D^n\times[0,1]$,
compactness gives a finite cubical subdivision whose closed boxes
each map into one trivializing neighborhood.
Process the time intervals in order, and the spatial boxes within
each interval in lexicographic order.
The prescribed part of a new box is a nonempty union of its lower
coordinate faces. After rescaling the box to $[0,1]^{n+1}$, this union
has a retraction
\[
 \rho_i(x)=
 \begin{cases}
  x_i-\min_{j\in S}x_j,&i\in S,\\
  x_i,&i\notin S,
 \end{cases}
\]
where $S$ indexes those faces.
In the chosen trivialization, extend the existing fibre coordinate
by composing it with $\rho$, and keep the prescribed base map.
The extensions agree on the already filled faces, and finite pasting
gives a continuous lift on the whole cube. This includes $n=0$.

(2) Apply Proposition~1.3.9 to the boundary inclusions and the map
$X\to\Delta^0$. By Exercise~\ref{ex:59}, this gives a monomorphism
$j\colon X\hookrightarrow W$ with $W\to\Delta^0$ a trivial fibration.
Thus $W$ is Kan and contractible by Exercise~\ref{ex:107}, and
$W\times Y\to Y$ is a trivial fibration by pullback.
The map $k=(j,p)\colon X\hookrightarrow W\times Y$ is a monomorphism.
Lifting against $p$ gives $r\colon W\times Y\to X$ such that
$rk=\mathrm{id}_X$ and $pr=\operatorname{pr}_Y$.
Hence $p$ is a retract of this product projection.
After realization, the projection is $|W|\times|Y|\to|Y|$,
which is Serre: a homotopy lifts by keeping the initial first coordinate.
The RLP passes to retracts, so $|p|$ is Serre as well.

Finally, use the minimal-fibration factorization stated before
Exercise~\ref{ex:107}. Both factors realize to Serre fibrations by
(1) and (2). Their composite is Serre, since one first lifts along
the second factor and then along the first. Thus $|p|$ is a Serre
fibration for every Kan fibration $p$, as asserted.
\end{proof}
\end{exercise}
